\section{Ascent, Staging and Upper Stage}
\label{sec:ascent}

This chapter describes the command laws behind the mission phases of
Section~\ref{sec:vehicle}: the ascent attitude program, the engine
shutdown, the separation state map with its audit, and the separated
upper-stage continuation. The event machine, step control and state bookkeeping of the run loop are
described in the architecture chapter of this report.

\subsection{Attitude program and ascent control}

Powered ascent flies a pitch program with an ideal attitude-control law.
The commanded tilt grows from vertical along
%
\begin{equation}
\label{eq:pitch-program}
\theta(t) = \theta_{\max}\,
\frac{x^2}{1+x^2},
\qquad x = \frac{t - t_s}{\tau},
\end{equation}
%
with $t_s = 8\,\mathrm{s}$, $\tau = 100\,\mathrm{s}$ and
$\theta_{\max} = 80^{\circ}$ toward a $45^{\circ}$ north-east heading (an
\ac{ISS}-launch approximation). The tilt law is calibrated against the
video-derived ascent reference used in the results chapter; it is not
independently validated guidance (class A). Attitude tracking applies a
proportional-derivative law on the Body-frame attitude error
$\mathbf{e}$ and angular rate,
%
\begin{equation}
\label{eq:attitude-pd}
\boldsymbol{\alpha} = k_p\,\mathbf{e} - k_d\,\boldsymbol{\omega},
\qquad
\mathbf{M} = \mathbf{J}\,\boldsymbol{\alpha},
\quad k_p = 2.0,\; k_d = 2.5,
\end{equation}
%
an unlimited ideal moment: the ascent controller relies on the
actuation chain only for realistic execution, and the provider optionally
enforces a declared angle-of-attack limit, raising a domain error when the
flight exceeds it.

\subsection{Engine shutdown and event guards}

Engine cutoff at $140\,\mathrm{s}$ follows a quintic flow envelope, not an
instantaneous stop:
%
\begin{equation}
\label{eq:quintic-shutdown}
\gamma(u) = (1-u)^3 (1 + 3u + 6u^2),
\qquad
\gamma'(u) = -\,\frac{30\,u^2 (1-u)^2}{T},
\quad u = \frac{t - t_0}{T},
\end{equation}
%
with $T = 1\,\mathrm{s}$ and completion anchored to the published
$141\,\mathrm{s}$ MECO reference. The envelope and its derivative are
supplied together to the integrator, keeping mass flow and flow rate
consistent during the ramp; the duration is a numerical design choice of
class A, anchored to, not measured from, the reference
timeline. Event transitions carry guards that must hold at the event:
separation requires zero flow, engines off and the full stack attached;
boostback, entry and landing ignitions additionally require the
respective propellant budget (including the reserve),
attitude readiness and propulsion readiness, and for the latter two a
descending trajectory. Each burn ends through a smooth ramp with its own
reason: predicted-miss tolerance, speed target, or budget cap.

\subsection{Stage separation and conservation audit}

Separation at $144\,\mathrm{s}$ is a zero-impulse, no-discard state map:
the coasting parent is split into booster and upper stack at the
stage-2 interface, both children keep the parent attitude and angular
rate, and the upper stage's stations are renumbered to its own aft-skirt
datum. The map refuses to run while any engine command or flow is active,
and its result is only returned if a conservation audit confirms
continuity to tight tolerances: mass, linear momentum, angular momentum
about the Earth-centred origin and kinetic energy are conserved to
$10^{-10}$ relative residual, fixed structural points keep their
ECEF position and velocity to $10^{-8}\,\mathrm{m/m/s}$, and attitude and
angular rate are preserved to $10^{-12}$. The audit quantities (all
residuals, and both child masses) are recorded in the run's event record,
and the results chapter quotes them.

\subsection{Upper-stage continuation}

After separation the upper stack continues under its own reduced provider
as a six-degree-of-freedom propagation from the conserved separation
state: a single MVac engine at $0.95$ throttle with quintic ignition and
cutoff ramps (the standard $10u^3 - 15u^4 + 6u^5$ envelope) over
$1\,\mathrm{s}$ at $152\,\mathrm{s}$ ignition and $545\,\mathrm{s}$
cutoff, an attitude law of the form of Equation~\eqref{eq:attitude-pd}
($k_p = 2.0$, $k_d = 2.5$) holding the separation heading, and a simple
axial-drag estimate of $C_D = 0.3$ on the $10.52\,\mathrm{m^2}$ stage
cross-section --- deliberately not the declared \texttt{second\_stack}
aerodynamic configuration of Section~\ref{sec:aero}. These are
demonstration commands of class A, not reconstructed flight data: the
continuation exists to show the separated state and the conserved
separation map, and its SECO state (about $200\,\mathrm{km}$ altitude at
$1\,600\,\mathrm{km}$ surface distance in the nominal run) feeds the
flight-path overview of Figure~\ref{fig:flightpaths} rather than any
orbital-accuracy claim.
